\section{Attack Surfaces in HPC}
\label{sec:taxonomy}

Prompt injection, poisoned tool output, and protocol abuse describe how an agent is manipulated. For HPC, we must also ask where the malicious content enters the workflow and what the redirected agent can affect. The following taxonomy is organized around those HPC surfaces. Its categories can lead to different harms, including data disclosure, allocation abuse, and corruption of scientific results, so they should not be treated as interchangeable instances of prompt injection.

\subsection{Shared parallel-filesystem poisoning}
\label{sec:taxonomy-filesystem}

Parallel filesystems and scratch spaces preserve state across jobs and, depending on site policy, expose shared or collaborative paths to several users. Malicious text placed in a result file, cache, README, or predictable output path may remain until a later agent reads it. The writer and reader need not run at the same time, which makes this a persistent workflow channel rather than a one-session interaction.

\subsection{Scheduler-induced co-location and log injection}
\label{sec:taxonomy-scheduler}

Agents that diagnose failures routinely inspect standard output, standard error, accounting records, and wrapper diagnostics. These inputs cannot simply be discarded as untrusted because they contain the evidence needed to repair a job. If an attacker can influence a related job or a shared output path, the resulting text enters a context in which scheduler commands are likely to be available. A successful redirection could therefore consume an allocation by submitting, resizing, or repeatedly resubmitting jobs.

\subsection{Tool, module, and MCP poisoning}
\label{sec:tool-poisoning}

HPC software stacks include site modules, group-maintained build helpers, post-processing scripts, and, increasingly, MCP tool servers. Trust is often based on a mixture of site policy and local convention. Shi et al.\ show that a crafted tool description, not only executable output, can bias an agent's tool selection \cite{Shi2025ToolHijacker}. In an HPC workflow, a module description or MCP manifest may therefore influence the agent before the selected program runs.

\subsection{Cross-project data leakage}
\label{sec:cross-project}

Researchers often belong to several projects, with access to separate allocations and datasets. An agent working on one project inherits whatever broader access the user exposes to it. A prompt-injection attack can exploit this mismatch by directing the agent to read from another project and disclose the result through an otherwise approved channel. Account-level permissions will not flag the read when the user is a legitimate member of both projects.

\subsection{Coordinated multi-agent exfiltration}
\label{sec:taxonomy-multiagent}

Multi-agent workflows pass messages and intermediate artifacts between planners and workers \cite{Rosendo2025ORNL,Mondesire2025PEARC}. If one agent or artifact is compromised, it can steer a second agent that has different access. The resulting disclosure may be spread across several ordinary-looking actions: one agent writes an artifact, another reads protected data, and a third sends context to an approved model backend. Detecting the sequence requires more than reviewing each agent's log in isolation.
